\documentclass[11pt,a4paper,twoside]{article}


\usepackage[margin=25mm, top=20mm]{geometry} % for setting page borders etc
\usepackage{amsmath}  % for math arrays and many math stuff
\usepackage{amsfonts} % for math fonts
\usepackage[colorlinks=true, allcolors=violet]{hyperref} % for hyperlink
\usepackage{multicol} % for multi-columns; \begin{multicols}{n}
\usepackage{bbding}
\usepackage{enumitem}
\usepackage[euler]{textgreek} % for greek letters in text

% \setlength{\parskip}{10pt}

%% Citations
\usepackage{nameref}
\usepackage{cleveref}
\crefformat{section}{\S#2#1#3}
\crefmultiformat{section}{\S\S#2#1#3}{ and~#2#1#3}{, #2#1#3}{, and~#2#1#3}
\crefformat{figure}{Fig.~#2#1#3}
% \citestyle{acmauthoryear}

%% ACM conference paper style font
\RequirePackage[T1]{fontenc}
\RequirePackage[tt=false, type1=true]{libertine}
\RequirePackage[varqu]{zi4}
\RequirePackage[libertine]{newtxmath}
\renewcommand\ttdefault{cmtt}

\usepackage{graphicx}
\usepackage[dvipsnames]{xcolor}

\usepackage{multirow}

\usepackage{scalerel}
\newcommand\shortarrow{\hstretch{0.7}{\rightarrow}}

\usepackage{listings} % for code listings; \begin{lstlisting}
\lstset{
  mathescape=true,
  frame=none,
  xleftmargin=0.8em,
  columns=fullflexible,
  basicstyle=\ttfamily,
  commentstyle={\color{gray!60!blue}\ttfamily},
  extendedchars=true,
  showstringspaces=false,
  aboveskip=\medskipamount,
  belowskip=\medskipamount,
  moredelim=**[is][\color{blue}]{/*}{*/},
  moredelim=**[is][\color{green!60!black}]{/!}{!/},
  moredelim=**[is][\color{gray}]{/(}{)/},
  moredelim=[is][\color{red}]{/[}{]/}
}

%%
\newtheorem{theorem}{Theorem}[section]
\newtheorem{lemma}[theorem]{Lemma}
\newtheorem{definition}[theorem]{Definition}
\newtheorem{coro}[theorem]{Corollary}

\newenvironment{syntax}{\[\begin{array}{l@{\quad}lcl}}{\end{array}\]}

\usepackage{bussproofs} % bussproofs

% Hack: wrap proof trees in a box as wide as the proof
% To fit multiple trees in one line
\newenvironment{bprooftree}
  {\leavevmode\hbox\bgroup}
  {\DisplayProof\egroup}

\usepackage{tikz} % for drawing diagrams
\usetikzlibrary{calc}
\usetikzlibrary{shapes}
\usetikzlibrary{positioning}
\usetikzlibrary{fit}
\usetikzlibrary{cd}

\raggedbottom

% Custom commands for typesetting
\newcommand{\semicolon}{\, ;}
\newcommand{\tightcolon}{{\:\!:\:\!}}
\newcommand{\gooddot}{.\;\!}

\newcommand{\lam}[3]{\lambda #1 \,{:}\, #2 .\; #3}
\newcommand{\Lam}[2]{\Lambda #1 \gooddot #2}
\newcommand{\app}[2]{#1 \, #2}
\newcommand{\fora}[2]{\forall #1 \gooddot #2}
\newcommand{\exist}[2]{\exists #1 \gooddot #2}
\newcommand{\newcon}{\cdot\, ; \cdot}
\newcommand{\var}[2]{#1 \,{:}\, #2}
\newcommand{\ifthen}[3]{\mathsf{if}\ #1 \ \mathsf{then}\ #2 \ \mathsf{else}\ #3}
\newcommand{\true}{\mathsf{true}}
\newcommand{\false}{\mathsf{false}}
\newcommand{\Bool}{\mathsf{Bool}}
\newcommand{\Halt}{\textbf{Halt}}
\newcommand{\sub}[2]{[#1] #2}
\newcommand{\pack}[3]{\mathsf{pack}_{#1}(#2,#3)}
\newcommand{\yes}{\mathsf{yes}}
\newcommand{\no}{\mathsf{no}}
\newcommand{\cho}{\mathsf{choose}}
\newcommand{\fst}{\mathsf{fst}}
\newcommand{\snd}{\mathsf{snd}}
\newcommand{\CList}{\mathsf{CList}}
\newcommand{\nat}{\mathsf{nat}}
\newcommand{\state}[2]{\langle #1, #2 \rangle}
\newcommand{\letseq}[3]{\mathsf{let}\ #1 = #2;\, #3}
\newcommand{\ret}[1]{\mathsf{return}\ #1}
\newcommand{\reduces}{\rightsquigarrow^*}
\newcommand{\type}{\mathsf{type}}
\newcommand{\sem}[1]{\lBrack #1 \rBrack}
\newcommand{\lamu}[3]{\mu #1 \,{:}\, #2 .\; #3}
\newcommand{\contra}[3]{\langle #1 \mid_{#2} #3 \rangle}
\newcommand{\refl}[1]{\mathsf{refl}\ #1}
\newcommand{\id}[3]{(#1 = #2 : #3)}
\newcommand{\subst}[5]{\mathsf{subst}[\var{#1}{#2} \gooddot #3](#4, #5)}
\newcommand{\Pitype}[3]{\Pi #1 \,{:}\, #2 \gooddot #3}
\newcommand{\Tcirc}[1]{{#1}^{\circ}}
\newcommand{\Tdot}[1]{{#1}^{\bullet}}
\newcommand{\qneg}[1]{\sim\!\!#1}
\newcommand{\Jrule}[4]{\mathsf{J}[#1 \gooddot #2](#3, #4)}
\newcommand{\hole}[1]{\mathcal{C}[#1]}


\begin{document}
\begin{center}
  {\Huge Part II Types: Example Sheet Solutions}\\
  \vspace{1.0em}
  {\large Yulong Huang and Yufeng Li}\\
  \vspace{0.5em}
  {\small Department of Computer Science and Technology \\ University
  of Cambridge \\ Oct 2026}\\
  \vspace{1.0em}
  \small{\textit{The tips and some questions are due to the previous
      example sheet by Nathanael Arkor, Dylan McDermott, Sam  Steenkamp,
      Domagoj Stolfa, et al., which has questions from Benjamin C.~Pierce
  and Andrew M.~Pitts.}}
\end{center}

\section*{Solutions}

This document provides solutions to the problems set in the three example sheets.
Answers to the simple warm-up questions are sometimes omitted or presented with
just the crucial steps.

\section*{Example Sheet 1}

\subsection*{I. Simply-Typed Lambda Calculus}

\paragraph{Question 1-6} Omitted. These questions are warn-up exercises to remind
the students what they learned in the Part IB Semantics course, and they should be 
able to answer these questions fairly easily. It is also for the supervisors to check
that students have correct understanding of the foundational concepts and good 
styles of written proofs.

\subsection*{II. Curry-Howard Correspondence}

\subsection*{Question 7}

\begin{enumerate}[label=(\alph*)]
  \item False is interpreted as the empty type. 
  \item Since there is no closed value of the empty type. So, by type-safety, any
  closed program of the empty type must loop forever. Therefore, termination
  proves that no closed program of the empty type exists, as it leads to a
  contradiction.
  \item No, a type system could be consistent while non-terminating. For example, we
    can extend STLC with the following construct
    \begin{center}
      \begin{bprooftree}
        \AxiomC{$\Gamma \vdash v : A$}
        \RightLabel{\scriptsize{Loop}}
        \UnaryInfC{$\Gamma \vdash \mathsf{loop}(v) : A$}
      \end{bprooftree}
      \qquad
      \begin{bprooftree}
        \AxiomC{}
        \UnaryInfC{$\mathsf{loop}(v) \leadsto \mathsf{loop}(v)$}
      \end{bprooftree}
    \end{center}
    which builds an infinite loop for every value of type $A$. Since there is no value of type $0$,
    it cannot be used to prove false in an empty context and consistency still holds.
\end{enumerate}
%
Students should be fully aware the of the motivation of proving termination: it is only required 
if a language is viewed as a logic. In reality almost all programming languages are not terminating
and termination does not matter. It is also important to know that this is a one-way implication.
Of course, in many cases ensuring termination is crucial, for example in dependently typed 
proof assisstants we must have logical consistency, and in eBPF we require termination
to safely execute code in the kernal.

\subsection*{Question 8} 

\begin{enumerate}[label=(\alph*)]
  \item The negation of $X$ is true if $X$ implies false.
  \item The first step is to translate the propositions into types, which is always
    in the form of $A_1 \to A_2 \to ... \to A_n$. Write down the outer level lambdas,
    and see if a term of type $A_n$ can be constructed with the bound variables of types
    $A_1\, ...\, A_{n-1}$.

    \begin{enumerate}[label=(\roman*)]
      \item $t_1 : (P \to 0) \times (Q \to 0) \to (P + Q) \to 0$\\
        $ t_1 = 
        \lam{f}{(P \to 0) \times (Q \to 0)}{
          \lam{x}{P + Q}{
            \case{x}{p}{q}{\app{(\fst\ f)}{p}}{\app{(\snd\ f)}{q}}
          }
        }$
      \item $t_2 : ((P + Q) \to 0) \to (P \to 0) \times (Q \to 0)$\\
        $ t_2 = 
        \lam{f}{(P + Q) \to 0}{
          (\lam{p}{P}{\app{f}{(\inl\ p)}}),
          (\lam{q}{Q}{\app{f}{(\inr\ q)}})}$
      \item $t_3 : (P \to 0) + (Q \to 0) \to P \times Q \to 0$\\
        $ t_3 = 
          \lam{f}{(P \to 0) + (Q \to 0)}{
            \lam{x}{P \times Q}{
              \case{f}{p}{q}{\app{p}{(\fst\ x)}}{\app{q}{(\snd\ x)}}
            }
          }$
      \item $t_4 : ((P \times Q) \to 0) \to (P \to 0) + (Q \to 0)$\\
        Impossible, cannot call the function of type $(P \times Q) \to 0$
        with just $p : P$ or $q : Q$.
    \end{enumerate}
  
  \item It means that, \emph{if interpreted directly}, STLC's corresponding logic proves
    less propositions than the propositional logic. Note that it is imprecise to say
    that STLC is weaker (proves less), as later in the course, we see that 
    show that propositional logic can be encoded into intuitionistic logic.
  
\end{enumerate}
%
A ``puzzle'' question to test the ability of type-driven programming.
It also prepares for the discussion of classical v.s. intuitionistic logic later.
In Example Sheet 3 Question 1(e), they can prove the final de Morgan rule with
classical logic.

\subsection*{III. Logical Relations}

\paragraph*{Question 9} 

\begin{enumerate}[label=(\alph*)]
  \item Omitted. Bookwork.
  \item No, the proof fails in the lambda case. Precisely, we want to show that
    if $\cdot \vdash \lam{x}{X}{e} : X \to Y$ then$ \lam{x}{X}{e} \in \Halt_{X \to Y}$.
    We need to show that $\sub{v / x}{e} \in \Halt_Y$ for all $v : X$, by definition,
    but the induction hypothesis cannot be applied to $\sub{v / x}{e}$.
  \item No, the similar problem appears again in the lambda case.
    We want to show that if $x : X \vdash \lam{y}{Y}{e} : Y \to Z$ and
    $v \in \Halt_X$ then $\sub{v / x}{(\lam{y}{Y}{e})} \in \Halt_{Y \to Z}$, which requires us to show
    that $\sub{v' / y}{\sub{v / x}{(\lam{y}{Y}{e})}}$ for all $v' \in \Halt_Y$.
    Again, we cannot apply the induction hypothesis in this case.

    The point of parts (b) and (c) is to make students re-invent the fundamental lemma.
    We cannot prove the desired theorem directly by induction; we need the technique of 
    \emph{proof by generalization}: showing a more general form such that 
    the desired statement is an instance of it. 
    We cannot generalize the statement as in part (c),
    and for similar reasons, we cannot generalize to cases with a fixed
    number of variables in the context. Therefore, we have to generalize it for an
    arbitrary context $\Gamma$ and all possible closing substitutions for $\Gamma$, 
    which leads to the form of the fundamental lemma.
  \item 
    We want to show that if $\Gamma \vdash \lam{x}{X}{e} : X \to Y$ and 
    $v_1, ..., v_n \in \Halt_{\Gamma}$, then 
      \[\sub{v_1 / x_1, ..., v_n / x_n}{(\lam{x}{X}{e})} \in \Halt_{X \to Y}.\]
    By definition, we need to show that for all $v \in \Halt_X$ we have
    $\sub{v_1 / x_1, ..., v_n / x_n, v / x}{e} \in \Halt_Y$, which is true by induction on
    the subderivation $\Gamma , x : X \vdash e : Y$.
  \item $e \in \Halt_{X \times Y}$ if
    (1) $e$ halts and (2) $e \leadsto^* (e_1, e_2)$ with $e_1 \in \Halt_X$ and $e_2 \in \Halt_Y$.
    Alternatively, it could be $e \in \Halt_{X \times Y}$
    if $(\fst\ e) \in \Halt_X$ and $(\snd\ e) \in \Halt_Y$.
    The two definitions should be equivalent, but the first one is recommended.

    The recipe of defining logical realtions for termination is that a term is in
    $\Halt$ if it reduces to a certain normal form and its sub-terms are also in $\Halt$
    \emph{in a logical way}, which means that it is expected to have be a logical conjunction
    for products, an implication for functions, etc. Therefore, the first form is
    recommended as it follows the recipe. The second form based on eliminators, 
    although equivalent, could fail if applied blindly to some other cases where
    the eliminator refers to some arbitrary type, for instance see the next question.
  \item $e \in \Halt_{X + Y}$ if (1) $e$ halts and
    (2)  $e \leadsto^* \inl\ e_1$ with $e_1 \in \Halt_X$ or $e \leadsto^* \inr\ e_2$
    with $e_2 \in \Halt_Y$.

    Some students are tempted to give the definition in terms of the eliminator of sums.
    For example, it could be like:
    $e \in \Halt_{X + Y}$ if $(\case{e}{x}{y}{e_1}{e_2}) \in \Halt_Z$ with certain conditions
    on $e_1$ and $e_2$. This is not a well-founded definition, as the $\Halt$ relation is
    defined \emph{inductively over the type}, this definition refers to an arbitrary type
    $Z$ which is not a sub-term of $X + Y$. The first definition follows the recipe given
    above and has no such problems.

\end{enumerate}

\subsection*{IV. System F}

\paragraph*{Question 10} 

\begin{enumerate}[label=(\alph*)]
  \item $\fora{\alpha}{\alpha} \to \fora{\beta}{\beta}$
  \item No solution, since $x : \alpha$ is not a big-lambda type.
  \item $\fora{\beta}{\beta}$
  \item $\fora{\beta}{\beta \to \beta}$
  \item No solution. $A_5$ must be $\alpha$ to match the typing judgement, but 
    that makes the function body ill-typed.
\end{enumerate}
%
Another puzzle question to familiarize with the rules of System F.

\section*{Example Sheet 2}

\subsection*{V. Church Encoding in System F}

\paragraph*{Question 1} 

\begin{enumerate}[label=(\alph*)]
  \item $\mathsf{unit} = \fora{\alpha}{\alpha \to \alpha}$\\
    $() = \Lam{\alpha}{\lam{x}{\alpha}{x}}$ 
  \item $\mathsf{empty} = \fora{\alpha}{\alpha}$\\
    $\mathsf{abort}_X\ e = \app{e}{X}$
  \item $\Bool = \fora{\alpha}{\alpha \to \alpha \to \alpha}$\\
    $\mathsf{true} = \Lam{\alpha}{\ulam{x}{\ulam{y}{x}}}$\\ 
    $\mathsf{false} = \Lam{\alpha}{\ulam{x}{\ulam{y}{y}}}$\\ 
    $\mathsf{if}_X\ e\ \mathsf{then}\ e_1\ \mathsf{else}\ {e_2} = e\ X\ e_1\ e_2$ 
  \item $\mathsf{nat} = \fora{\alpha}{\alpha \to (\alpha \to \alpha) \to \alpha}$\\
    $\mathsf{zero} = \Lam{\alpha}{\ulam{z}{\ulam{s}{z}}}$\\ 
    $\mathsf{suc}\ e = \Lam{\alpha}{\ulam{z}{\ulam{s}{\app{s}{(e\ \alpha\ z\ s)}}}}$\\ 
    $\mathsf{iter}_X(e,\ 0 \mapsto e_z,\ \mathsf{suc}\ z \mapsto e_s) = e\ X\ e_z\ e_s$    
  \item $\mathsf{option} A = \fora{\alpha}{\alpha \to (A \to \alpha) \to \alpha}$\\
    $\mathsf{none} = \Lam{\alpha}{\ulam{z}{\ulam{f}{z}}}$\\ 
    $\mathsf{some}\ e = \Lam{\alpha}{\ulam{z}{\ulam{f}{\app{f}{e}}}}$\\ 
    $\mathsf{case}_X(e,\ \mathsf{none} \mapsto e_n,\ \mathsf{some}\ z \mapsto e_s) = e\ X\ e_n\ e_s$    
\end{enumerate}
%
Church encoding is the most examed topic and very easy to learn with some practice.
Students should be able to spot the pattern from completing this question: 
the church encoding is defined in the form of
$\fora{\alpha}{A_1 \to ... \to A_n \to \alpha}$, where each constructor of the datatype 
corresponds to an argument $A_i$, and each $A_i$ is in the form of $B_1 \to ... \to B_m \to \alpha$,
where $B_i$ corresponds to the arguments that constructor takes. 
It is then straightforward to build the church-encoded constructors and the eliminators
are simply applications.

\paragraph*{Question 2}

\begin{enumerate}[label=(\alph*)]
  \item $\mathsf{pred} = \ulam{n}{\fst\ (\mathsf{iter}\ \mathsf{nat}\ (0,0)\ (\ulam{r}{(\snd\ r, \mathsf{suc}\ (\snd\ r))})\ n)}$.
  
  The key is to construct the following sequence
    \[ (0,0), (0,1), (1,2), (2,3), ...\]
    where the $n$-th term can be computed with 
    $(\mathsf{iter}\ \mathsf{nat}\ (0,0) (\ulam{r}{(\snd\ r, \mathsf{suc}\ (\snd\ r))})\ n)$.
    Taking the first projection of this expression gives the predecessor.
  \item The key is to construct the following sequence
    \[ \mathsf{none}, \mathsf{some}\ 0, \mathsf{some}\ 1, \mathsf{some}\ 2, ...\]
    using $\mathsf{none}$ and $\mathsf{some}$ to distinguish the zero and succesor case
    of the predecessor's input. The $n$-th term of the sequence can be computed with
    \[
      f\ n = \mathsf{iter}\ (\mathsf{option}\ \mathsf{nat})\ \mathsf{none}\ 
      (\ulam{r}{\mathsf{case}(r,\ \mathsf{none} \mapsto \mathsf{some}\ 0,\ 
      \mathsf{some}\ z \mapsto \mathsf{some}\ (\mathsf{suc}\ z))}).\]
    We then define a function that maps the sequence back to natural numbers,
    sending $\mathsf{none}$ to $0$ and $\mathsf{suc}\ z$ to $z$:
    \[ 
      g\ x = \mathsf{case}(x,\ \mathsf{none} \mapsto 0,\ \mathsf{some}\ z \mapsto z)
      \]
    and the predecessor of $n$ is simply $g\ (f\ n)$.
    
  \item O(n) steps.
\end{enumerate}
%
Predecessor are also likely to appear in exams.
Without the hints, this question would be very difficult to solve, but
the hints and the warm-up in the previous question prepares students well enough
for the challange. Those who understands the question and knows how to apply it 
will have no problems with church encodings in exams.

Sharper students might notice that, fundamentally, option types are used here
because of the isomorphism $\mathbb{N} \cong 1 + \mathbb{N}$, and 
the similar technique could be used to define pred-like functions for other datatypes,
such as head and tail for lists. They might also relate the iterator with the generalized
recursor in Example Sheet 3 Question 9.

\paragraph*{Question 3}
\begin{enumerate}[label=(\alph*)]
  \item $\mathsf{fmap} = 
    \Lam{\beta_1}{\Lam{\beta_2}{
      \lam{f}{(\beta_1 \to \beta_2)}{
        \lam{x}{(1+\beta_1)}{
          \case{x}{()}{y}{\inl\ ()}{\inr\ (f\ y)}
        }}}}$
  \item $\mathsf{fold} = 
    \Lam{\alpha}{
      \lam{f}{(1 + \alpha \to \alpha)}{
        \lam{n}{(\fora{\beta}{(1+\beta \to \beta) \to \beta})}{
          n\ \alpha\ f
        }}}$
  \item $\mathsf{zero} = 
    \Lam{\beta}{
      \lam{f}{(1+\beta \to \beta)}{f\ (\inl\ ())}
    }$\\
    $\mathsf{suc} = 
    \lam{n}{(\fora{\beta}{(1+\beta \to \beta) \to \beta})}
      {\Lam{\beta}{
        \lam{f}{(1+\beta \to \beta)}{f\ (\inr\ (n\ \beta\ f))}
      }}$\\
    Given an eliminator $\mathsf{iter}_X(e,\ 0 \mapsto e_z,\ \mathsf{suc}\ z \mapsto e_s)$,
    define 
      \[f = \lam{x}{(1 + X)}{\case{x}{()}{y}{e_z}{e_s\ y}}.\]
    Then, the iterator is encoded as $(\mathsf{fold}\ X\ f\ n)$, which $\beta$-reduces to
    $(n\ X\ f)$. The encoding has the desired behaviour:
    \begin{itemize}
      \item $\mathsf{iter}_X(\mathsf{zero},\ 0 \mapsto e_z,\ \mathsf{suc}\ z \mapsto e_s)
        = \mathsf{zero}\ X\ f = f\ (\inl\ ()) = e_z$
      \item $\mathsf{iter}_X(\mathsf{suc}\ n,\ 0 \mapsto e_z,\ \mathsf{suc}\ z \mapsto e_s)
        = (\mathsf{suc}\ n)\ X\ f = f\ (\inr\ (n\ X\ f)) = e_s\ (n\ X\ f)
        = e_s\ \mathsf{iter}_X(n,\ 0 \mapsto e_z,\ \mathsf{suc}\ z \mapsto e_s)$
    \end{itemize}
    So, $\mu F$ is indeed a valid encoding of the natural number.

\end{enumerate}
%
This question gives students a taste of algebraic datatypes as initial algebras
of polynomial functors. Question B2 considers the more general case. It is not supposed
to be a hard question, but one that makes the students realize that church encoding
is not an ad-hoc process driven by intuition but a sturctured, systemetic technique.

\paragraph*{Question 4}
\begin{enumerate}[label=(\alph*)]
  \item $\mathsf{and} = \ulam{b_1}{\ulam{b_2}{b_1\ \mathsf{CBool}\ b_2\ b_1}}$
  \item $\mathsf{or} = \ulam{b_1}{\ulam{b_2}{b_1\ \mathsf{CBool}\ b_1\ b_2}}$
\end{enumerate}
%
The equivalent version of $a \land b = \ifthen{a}{b}{a}$ and 
$a \lor b = \ifthen{a}{a}{b}$. Warm-up for the next exercise.

\subsection*{VI. Existential and Semantic Types}

\paragraph*{Question 5}
\begin{enumerate}[label=(\alph*)]
  \item $\mathsf{and} = \ulam{b_1}{\ulam{b_2}{\cho\ \beta\ b_1\ b_2\ b_1}} $
  \item $\mathsf{extend} = \lam{b}{\exist{\beta}{\Bool}}{
    \mathsf{let\ pack}\ (\beta, bool) = b\  \mathsf{in}\ 
    \pack{\Bool \times \beta \to \beta \to \beta}{\beta}{(bool, \mathsf{and})}} $
  \item $\mathsf{Nat} = \beta \times (\beta \to \beta) \times 
    \fora{\alpha}{\beta \to \alpha \to (\alpha \to \alpha) \to \alpha} $
\end{enumerate}
%
An important question to make students familiar with existential types and using them as 
packages.

\paragraph*{Question 6}
\begin{enumerate}[label=(\alph*)]
  \item Proof goal: if $e \leadsto e'$ and $e \in \sem{\Theta \vdash \fora{\alpha}{A}\ \type}\ \theta$,
    then $e' \in \sem{\Theta \vdash \fora{\alpha}{A}\ \type}\ \theta$.
    It suffices to show that (1) $e'$ halts and (2) $\forall (A', X).\ 
    e'\ A' \in \sem{\Theta, \alpha \vdash A\ \type}\ (\theta, X / \alpha)$. 
    We have (1) since $e \in \sem{\Theta \vdash \fora{\alpha}{A}\ \type}\ \theta$, so $e$ halts,
    and therefore $e'$ halts because $e \leadsto e'$.
    We also know that $e\ A' \leadsto e'\ A'$ and $e\ A' \in \sem{\Theta, \alpha \vdash A\ \type}\ (\theta, X / \alpha)$.
    So, we can prove (2) by I.H.
  \item Proof goal: if $e \in \sem{\Theta \vdash \sub{A / \alpha}{(\fora{\beta}{B})}\ \type}\ \theta$,
    then $e \in \sem{\Theta , \alpha \vdash \fora{\beta}{B}\ \type}\ (\theta, \sem{\Theta \vdash A\ \type}\ \theta / \alpha)$.
    It suffices to show that $\forall (B, Y).\ e\ B \in \sem{\Theta, \alpha, \beta \vdash B\ \type}\ (\theta, \sem{\Theta \vdash A\ \type}\ \theta / \alpha, Y / \beta)$.
    Indeed, from assumption, we have
    \begin{align*}
      & e \in \sem{\Theta \vdash \fora{\beta}{\sub{A / \alpha}{B}}\ \type}\ \theta\\
      \Longrightarrow\ & \forall (B, Y).\ e\ B \in \sem{\Theta, \beta \vdash \sub{A / \alpha}{B}\ \type}\ (\theta, Y / \beta) \quad \text{by def.}\\
      \Longrightarrow\ & \forall (B, Y).\ e\ B \in \sem{\Theta, \beta, \alpha \vdash {B}\ \type}\ (\theta, Y / \beta, \sem{\Theta, \beta \vdash A\ \type}\ (\theta,  Y / \beta) / \alpha) \quad \text{by I.H.}\\
      \Longrightarrow\ & \forall (B, Y).\ e\ B \in \sem{\Theta, \beta, \alpha \vdash {B}\ \type}\ (\theta, Y / \beta, \sem{\Theta \vdash A\ \type}\ \theta / \alpha) \quad \text{by weakening}\\
      \Longrightarrow\ & \forall (B, Y).\ e\ B \in \sem{\Theta, \alpha, \beta \vdash B\ \type}\ (\theta, \sem{\Theta \vdash A\ \type}\ \theta / \alpha, Y / \beta) \quad \text{by exchange.}
    \end{align*}
  \item Proof goal: if $\Theta; \Gamma \vdash \Lam{\alpha}{e} : \fora{\alpha}{B}$, $\Theta \vdash \Gamma\ \mathsf{ctx}$,
    $\theta$ interprets $\Theta$, $e_i \in \sem{\Theta \vdash A_i\ \type}\ \theta$, 
    then $\sub{\overrightarrow{e} / \Gamma}{\sub{\overrightarrow{C} / \theta}{(\Lam{\alpha}{e})}} 
    \in \sem{\Theta \vdash \fora{\alpha}{B}\ \type}\ \theta$.
    Let $e' = \sub{\overrightarrow{e} / \Gamma}{\sub{\overrightarrow{C} / \theta}{{e}}}$.
    It suffices to show that $\forall (C, X).\ (\Lam{\alpha}{e'})\ C \in \sem{\Theta , \alpha \vdash B\ \type}\ (\theta, X / \alpha)$.
    We have $(\Lam{\alpha}{e'})\ C \leadsto \sub{C / \alpha}{e'}$. 
    We can apply I.H. to show that $\sub{C / \alpha}{e'} \in \sem{\Theta , \alpha \vdash B\ \type}\ (\theta, X / \alpha)$
    and prove the goal by closure. All conditions required to apply the I.H. can be easily derived from the assumption.
  \end{enumerate}
%
Semantic types is way too technical and rarely examed. This question could be skipped
if there is no time for it.


\subsection*{VII. Type Systems with States}

\paragraph*{Question 7}
\begin{enumerate}[label=(\alph*)]
  \item Omitted. Bookwork.
  \item The extended rules are given as follows. There could be other ways to give types and 
    operational semantics, for example, freeing a location could return the value stored in that location.
    \begin{center}
    \begin{bprooftree}
      \AxiomC{$\Sigma; \Gamma \vdash e : \mathsf{ref}\ X$}
      \RightLabel{\scriptsize{Free}}
      \UnaryInfC{$\Sigma; \Gamma \vdash \mathsf{free}(e) : 1$}
    \end{bprooftree}
    \qquad
    \begin{bprooftree}
      \AxiomC{}
      \UnaryInfC{$\langle \sigma, l : v \semicolon \mathsf{free}(l) \rangle \leadsto 
      \langle \sigma \semicolon () \rangle$}
    \end{bprooftree}
    \qquad
    \begin{bprooftree}
      \AxiomC{$\langle \sigma, e \rangle \leadsto \langle \sigma', e' \rangle$}
      \UnaryInfC{$\langle \sigma, \mathsf{free}(e) \rangle \leadsto \langle \sigma', \mathsf{free}(e') \rangle$}
    \end{bprooftree}
  \end{center}
  \item Not anymore, since double-free and use-after-free are possible.
\end{enumerate}
%
A simple question to test the understanding of type safely with states.

\paragraph*{Question 8}
\begin{enumerate}[label=(\alph*)]
  \item \lstinline|f : (int -> int) -> int -> int|\\
    \lstinline|r : (int -> int) ref|\\
    \lstinline|recur : int -> int|
  \item \lstinline|knot f = recur x = !r x = f recur x|
  \item $\mathsf{fib} = $ \lstinline|knot (fun r n -> if n = 0 or n = 1 then 1 else r (n - 1) + r (n - 2))|\\
    $\mathsf{loop} = $ \lstinline|knot (fun r -> r)|
  \item Initialize the reference with \lstinline|r = ref (fun x -> x)| will do.
\end{enumerate}
%
The point of this question is to make explicit connection that Landin's knot
implements a fixed-point operator, and make sure that students know how to program
with fixed points, in case they did not choose the Denotational Semantics course.
Also helps to memorize the knot's construction.

\subsection*{VIII. Monadic Type Systems}

\paragraph*{Question 9}

\begin{align*}
  \mathsf{knot} = \ulam{f}{ \{
    \mathsf{let}\ (r : \mathsf{ref}(\mathsf{int} \to T\ \mathsf{int}))  &= \{\mathsf{new}\ (\ulam{x}{\{\mathsf{return}\ 0\}})\}; \\
    \mathsf{let}\ (recur : \mathsf{int} \to T\ \mathsf{int}) &= \{ \mathsf{return}\ (\ulam{n}{
      \{ \letseq{r'}{\{ !r \}}{\mathsf{return}\ (f\ r'\ n)}\}}) \}; \\
    \mathsf{let}\ (\_ : 1) &= \{ r := recur \}; \\
    & \mathsf{return}\ recur \}
  }
\end{align*}
%
It is very easy to make small mistakes in this question. 
I have not implemented a type-checker for the monadic language so there might be
typo or mistakes in the code above as well.
The purpose is not to be fully correct but to make students write some complicated programs with the monadic language,
as its syntax and typing rules are not very intuitive.

\paragraph*{Question 10}
\begin{enumerate}[label=(\alph*)]
  \item
    \begin{enumerate}[label=(\roman*)]
      \item $\eta_X = \ulam{x}{\{\mathsf{return}\ x\}}$
      \item $\mu_X = \ulam{t}{\{\letseq{x}{t}{\letseq{y}{x}{y}}\}}$
      \item $\mathsf{fmap} = \ulam{f}{\ulam{t}{\{\letseq{x}{t}{\mathsf{return}\ (f\ x)}\}}}$ 
    \end{enumerate}
  \item    
    \begin{enumerate}[label=(\roman*)]
      \item \begin{align*}
          & \state{\sigma}{\letseq{y}{\mu_X(\app{\eta_{TX}}{v})}{\ret{y}}} \\
          \reduces & \state{\sigma}{\letseq{y}{\mu_X(\{\ret{v}\})}{\ret{y}}} \\
          \reduces & \state{\sigma}{\letseq{y}{\letseq{z}{v}{z}}{\ret{y}}} \\
          \reduces & \state{\sigma}{\letseq{y}{\ret{v'}}{\ret{y}}} \\
          \reduces & \state{\sigma}{\ret{v'}}
        \end{align*}
      \item \begin{align*}
          & \state{\sigma}{\letseq{y}{\app{\app{\mathsf{fmap}}{f}}{\{\ret{v}\}}}{\ret{y}}}\\
          \reduces & \state{\sigma}{\letseq{y}{\ret{f\ v}}{\ret{y}}} \\
          \reduces & \state{\sigma}{\ret{v'}}
        \end{align*}
    \end{enumerate}
  \item $\mathsf{bind} = \ulam{t}{\ulam{f}{\{\letseq{x}{t}{\letseq{y}{f\ x}{\ret{y}}} \}}}$.
  Alternatively, $\mathsf{bind}\ x\ f = \mu_{TY}(\mathsf{fmap}\ f\ x)$.
\end{enumerate}
%
This is another category-theoretic question, connecting the computational monad
(with return and bind) with the categorical monad (with $\mu$ and $\eta$).
The student might be unfamiliar with both, and in that case,
covering the computational notion is crucial (more important than
the monadic language taught in this course!) and the categorical definition
is to spark their interests.

\section*{Example Sheet 3}

\subsection*{IX. Classical Logic}

\paragraph*{Question 1} Omitted. Warm-up. Note that part (e) is done by a contradiction
with $\neg (A \land B)$ followed by a contradiction with $\neg A \lor \neg B$.

\paragraph*{Question 2} Omitted. Warm-up. 

\paragraph*{Question 3}
\begin{enumerate}[label=(\alph*)]
  \item Omitted.
  \item $\contra{\lamu{u}{A}{\contra{p}{A}{u}}}{A}{k} \mapsto \sub{k / u}{\contra{p}{A}{u}}
  \mapsto \contra{p}{A}{k}$.
  \item $\mathsf{upc}$ stands for \emph{unnecessary proof by contradiction}, which 
    proves $A$ by assuming that $A$ is false, then proves $A$, and then prove $A$
    by contradiction. 
\end{enumerate}
%
A small question to test the understanding of Curry-Howard correspondence
in this calculus. Occasionally students write unintentional UPCs and this
is a gentle reminder!

\subsection*{X. Continuation-Passing Style}

\paragraph*{Question 4}
\begin{enumerate}[label=(\alph*)]
  \item $\lam{k}{\neg \neg X}{\letcont{u}{\neg X}{\throw{X}{(u,k)}}}$
  \item $\letcont{u}{\neg(X \lor \neg X)}{\throw{X \lor \neg X}{
    (u, \inr\ (\letcont{v}{\neg \neg X}{\throw{\neg X}{(u, \inl\ (\mathsf{dne}\ v))}}))
    }}$\\
    Alternatively, 
    $\letcont{u}{\neg(X \lor \neg X)}{\throw{X \lor \neg X}{
      (u, \inl\ (\letcont{v}{\neg X}{\throw{X}{(u, \inr\ v)}}))
    }}$
  \item $\lam{f}{((X \to Y) \to X)}{\letcont{u}{\neg X}{f\ (\lam{x}{X}{\throw{Y}{(u, x)}})}}$
  \item $\lam{f}{(X \to Y)}{\lam{u}{\neg Y}{\letcont{v}{\neg\neg X}{
    \throw{\neg X}{(u, f\ (\mathsf{dne}\ v))}}}}$
  \item $\lam{k}{\neg(A \land B)}{\letcont{u}{\neg(\neg A \lor \neg B)}{
    \throw{\neg A \lor \neg B}{(k, (e_1, e_2))}
  }}$\\
    where $e_1 = \letcont{v}{\neg A}{\throw{A}{(u, \inl\ v)}}$,
    $e_2 = \letcont{v}{\neg B}{\throw{B}{(u, \inr\ v)}}$
\end{enumerate}

\paragraph*{Question 5}
\begin{enumerate}[label=(\alph*)]
  \item We create a program point that acts like a ``trap'' and pass it to $f$.
    Inside $f$, if it uses its argument by applying it to some term of type $X$, 
    the trap is activated and captures the value of type $X$ and returns it to us. 
    Otherwise, $f$ produces some term of type $X$ on its own. 
    Either way, we eventually get a term of type $X$.
  \item $\lam{xs}{\mathsf{List}\ X}
    \letcont{u}{\neg\nat}{\mathsf{fold}(xs, \mathsf{[]} \mapsto 0, (x, r) \mapsto \ifthen{x = 0}{\throw{\nat}{(u, 0)}}{x * r})}$ 
\end{enumerate}

\paragraph*{Question 6}
\begin{enumerate}[label=(\alph*)]
  \item $\Tcirc{1} = 1$, $\Tcirc{\langle\rangle} = \lam{k}{\qneg{1}}{\app{k}{\langle\rangle}} \,:\ \sim\qneg{1}$.
    The transformation turns a value into a function that given any continuation,
    supply the value into that continuation.
  \item $\Tcirc{X \to Y} = \Tdot{X} \to \Tdot{Y}$\\
    $\Tdot{(\lam{x}{X}{e})} = \lam{k}{\qneg(\Tdot{X} \to \Tdot{Y})}{\app{k}{(\lam{x}{\Tdot{X}}{\Tdot{e}})}} \,:\ \sim\qneg{(\Tdot{X} \to \Tdot{Y})}$\\
    $\Tdot{(\app{e_1}{e_2})} = \lam{k}{\qneg{\Tcirc{Y}}}{\Tdot{e_1}\ (\lam{f}{\Tdot{X} \to \Tdot{Y}}{f\ \Tdot{e_2}\ k})} \,:\ \sim\qneg{\Tcirc{Y}}$ \\

    $\Tcirc{(X \times Y)} = \Tdot{X} \times \Tdot{Y} $\\
    $\Tdot{(e_1, e_2)} = \lam{k}{\qneg{\Tdot{X} \times \Tdot{Y}}}{k\ (\Tdot{e_1}, \Tdot{e_2})} \,:\ \sim\qneg{(\Tdot{X} \times \Tdot{Y})}$
  \item $\Tcirc{(\neg{X})} =\ \qneg{\Tcirc{X}}$ \\
    $\Tdot{(\letcont{u}{\neg X}{e})} = \lam{k}{\qneg{\Tcirc{X}}}{(\lam{u}{\sim\sim\qneg{\Tcirc{X}}}{\Tdot{e}})\ (\mathsf{dni}(k))}$
    where $\mathsf{dni} = \lam{x}{X}{\lam{k}{\qneg{X}}}{k\ x}$. The inner application could also be written as a substitution,
    as in the lecture slides.
  \item Omitted. This should be done through parts (a-c).
\end{enumerate}

\paragraph*{Question 7} 
\begin{align*}
  \textsf{eq-at}\ xs\ ys & = \\
  &\letseq{i}{\mathsf{amb}\ (\mathsf{length}\ xs)}{} \\
  &\letseq{j}{\mathsf{amb}\ (\mathsf{length}\ ys)}{} \\
  &\mathsf{assert}(\mathsf{nth}\ i\ xs = \mathsf{nth\ j\ ys}); \\
  &(i, j)
\end{align*}

\subsection*{XI. Dependent Types in Programming}

\paragraph*{Question 8}
\begin{enumerate}[label=(\alph*)]
  \item \lstinline|{A : Set}| means that \lstinline|A| is a type (\lstinline|Set|
    is the type of small types) and the braces means that it is an implicit argument
    to the function. Implicit arguments are typically inferred, since if we know the 
    type of the input \lstinline|List Nat| we can infer that \lstinline|A = Nat|.
    Implicit arguments can also be supplied explicitly with syntax \lstinline|hd {A = A} ls|.
  \item The definition is given below.
    \begin{lstlisting}
      data Vec (A : Set) : (n : Nat) -> Set where
        [] : Vec A 0
        _::_ : {n : Nat} -> A -> Vec A n -> Vec A (suc n)
    \end{lstlisting}
  \item 
    \begin{enumerate}[label=(\roman*)]
      \item See code below.
        \begin{lstlisting}
          hd : {n : Nat} -> Vec A (suc n) -> A
          hd (x :: xs) = x

          tl : {n : Nat} -> Vec A (suc n) -> Vec A n
          tl (x :: xs) = xs

          append : {m n : Nat} -> Vec A m -> Vec A n -> Vec A (m + n)
          append [] ys = ys 
          append (x :: xs) ys = x :: (append xs ys)
        \end{lstlisting}
      \item See code below.
        \begin{lstlisting}
          zip : {n : Nat} -> Vec A n -> Vec B n -> Vec (A x B) n
          zip [] [] = []
          zip (x :: xs) (y :: ys) = (x , y) :: zip xs ys
        \end{lstlisting}
    \end{enumerate}
    These are total functions and their counterparts for lists are partial functions.
  \item No, there is no way to supply a term of type \lstinline|A| for the empty case.
    In general, we cannot define partial functions in Agda (directly, we could for example
    encode with the maybe monad).
\end{enumerate}

\subsection*{XII. Dependent Types in Theorem Proving}

\paragraph*{Question 9}
\begin{enumerate}[label=(\alph*)]
  \item It corresponds to the $\forall$ operator: given an element $a \in A$, 
    we get a proof of the predicate $B(a)$.
  \item The code is given as the follows.
  \begin{lstlisting}
    idr : (n : Nat) -> add n Z == n
    idr Z = refl
    idr (S n) = cong S (idr n)
  \end{lstlisting}
  It is a proof by induction. We cannot use \lstinline|refl| directly in the succesor
  case because \lstinline|add (S n) Z| is not \emph{definitionally} equal to \lstinline|S n|.
  More specifically, the left hand side cannot convert to the right hand side
  by the reflexive-transitive closure of reductions (this is an over-simplified explaination
  that is not always applicable, since conversions can be defined without reductions).
  \item The induction principle for the natural numbers.
  \item The generalized iterator, where the successor's case has access
    to the predecessor. This is also called a recursor.
\end{enumerate}
%
This question gives students a taste of proof by induction in Agda.
Being able to read \lstinline|elim| write \lstinline|idr| is enough for the exam. 
The more important thing is to see the Curry-Howard correspondence 
with dependent types.

\paragraph*{Question 10}
\begin{enumerate}[label=(\alph*)]
  \item Omitted. Bookwork.
  \item $\mathsf{trans}_A = \ulam{x}{\ulam{y}{\ulam{z}{
    \lam{p}{\id{x}{y}{A}}{
      \lam{q}{\id{y}{z}{A}}{
        \subst{w}{A}{\id{x}{w}{A}}{q}{p}
      }}}}}$
  \item $\mathsf{cong}_{A,B} = \ulam{x}{\ulam{y}{\ulam{f}{
    \lam{p}{\id{x}{y}{A}}{
        \subst{w}{A}{\id{f\ x}{f\ w}{B}}{p}{\refl{(f\ x)}}
      }}}}$
\end{enumerate}
%
This question shows students some fundamental equational reasoning with
the identity type and its (simplified) elimination rule, which could be
not intuitive at first. The full elimination rule is explored further
in Question C3.

\end{document}
